\documentclass[11pt,a4paper]{article}

% ==============================================================================
% PAKIETY JĘZYKOWE, KODOWANIE I TYPOGRAFIA
% ==============================================================================
\usepackage[utf8]{inputenc}
\usepackage[T1]{fontenc}
\usepackage[polish]{babel}
\usepackage{lmodern}

% Poprawa czytelności i eliminacja pojedynczych wiszących linii (wdów i sierot)
\usepackage{microtype}
\clubpenalty=10000
\widowpenalty=10000
\displaywidowpenalty=10000

% ==============================================================================
% PAKIETY MATEMATYCZNE
% ==============================================================================
\usepackage{amsmath, amssymb, amsfonts, amsthm}
\usepackage{mathtools}

% ==============================================================================
% UKŁAD STRONY I ODSTĘPY
% ==============================================================================
\usepackage[margin=2.3cm, top=2.5cm, bottom=2.5cm, headheight=14.5pt]{geometry}
\usepackage{fancyhdr}
\usepackage{lastpage}
\usepackage{enumitem}
\usepackage{booktabs}
\usepackage{array}
\usepackage{graphicx}
\usepackage{xcolor}

% Ustawienia list wyliczeniowych
\setlist{itemsep=3pt, topsep=4pt, parsep=1pt}

% ==============================================================================
% PALETA KOLORÓW I BLOKI TCOLORBOX
% ==============================================================================
\definecolor{navyblue}{HTML}{1A365D}     % Główny akcent, nagłówki
\definecolor{forestgreen}{HTML}{1E6B52}  % Kolor sekcji teorii
\definecolor{charcoal}{HTML}{2D3748}     % Kolor podsumowań
\definecolor{lightbg}{HTML}{F7FAFC}      % Jasne tło treści zadania
\definecolor{teoriabg}{HTML}{F2F8F5}     % Jasne tło teorii
\definecolor{darkslate}{HTML}{4A5568}    % Tekst poboczny

\usepackage{tcolorbox}
\tcbuselibrary{skins,breakable}

% Konfiguracja boksów z nowoczesnym paskiem bocznym (Bourbaki / AMS style)
\newtcolorbox{tresczadania}[1][]{
    enhanced,
    breakable,
    colback=lightbg,
    colframe=navyblue,
    borderline west={3.5pt}{0pt}{navyblue},
    boxrule=0pt,
    arc=0mm,
    top=7pt, bottom=7pt, left=10pt, right=10pt,
    fonttitle=\bfseries\color{navyblue},
    title={Treść zadania},
    attach title to upper={\par\vspace{5pt}},
    #1
}

\newtcolorbox{teoriabox}[1][Wymagana teoria i najważniejsze wzory]{
    enhanced,
    breakable,
    colback=teoriabg,
    colframe=forestgreen,
    borderline west={3.5pt}{0pt}{forestgreen},
    boxrule=0pt,
    arc=0mm,
    top=7pt, bottom=7pt, left=10pt, right=10pt,
    fonttitle=\bfseries\color{forestgreen},
    title={#1},
    attach title to upper={\par\vspace{5pt}}
}

\newtcolorbox{odpowiedzbox}[1][Podsumowanie i wynik]{
    enhanced,
    breakable,
    colback=gray!6!white,
    colframe=charcoal,
    borderline west={3.5pt}{0pt}{charcoal},
    boxrule=0pt,
    arc=0mm,
    top=7pt, bottom=7pt, left=10pt, right=10pt,
    fonttitle=\bfseries\color{charcoal},
    title={#1},
    attach title to upper={\par\vspace{5pt}}
}

% ==============================================================================
% KONFIGURACJA NAGŁÓWKÓW I STOPEK
% ==============================================================================
\pagestyle{fancy}
\fancyhf{}
\fancyhead[L]{\small\textcolor{darkslate}{\textsc{Matematyka 1} \textbullet\ Lista 6: Szeregi liczbowe}}
\fancyhead[R]{\small\textcolor{darkslate}{\textsc{Rozwiązania zadań}}}
\fancyfoot[C]{\small\textcolor{darkslate}{Strona \thepage\ z \pageref{LastPage}}}
\renewcommand{\headrulewidth}{0.4pt}
\renewcommand{\footrulewidth}{0.4pt}

% ==============================================================================
% HIPERŁĄCZA (HYPERREF)
% ==============================================================================
\usepackage{hyperref}
\hypersetup{
    colorlinks=true,
    linkcolor=navyblue,
    urlcolor=navyblue,
    citecolor=forestgreen,
    pdftitle={Matematyka 1 - Lista 6: Szeregi liczbowe (Rozwiązania)},
    pdfauthor={Rozwiązania zadań},
    pdfsubject={Szeregi liczbowe, kryteria zbieżności, dwumian Newtona}
}

% ==============================================================================
% POCZĄTEK DOKUMENTU
% ==============================================================================
\begin{document}

% ------------------------------------------------------------------------------
% NAGŁÓWEK DOKUMENTU (ACADEMIC HEADER BANNER)
% ------------------------------------------------------------------------------
\begin{center}
    {\scshape\Large Adam Gwiazdoń}\\[0.2cm]
    {\scshape\large Kurs: Matematyka 1 --- Semestr Zimowy}\\[0.4cm]
    \rule{\textwidth}{1.2pt}\\[0.35cm]
    {\huge\bfseries\color{navyblue} Rozwiązania zadań: Lista 6}\\[0.2cm]
    {\large\bfseries Tematyka: Dwumian Newtona oraz kryteria zbieżności szeregów liczbowych}\\[0.3cm]
    \rule{\textwidth}{0.8pt}\\[0.3cm]
    {\small \textbf{Opracowanie:} Praca domowa \quad\textbullet\quad \textbf{Data:} \today}
\end{center}

\vspace{0.4cm}

\begin{abstract}
\noindent Niniejszy dokument zawiera kompletne, formalne i szczegółowe rozwiązania wszystkich zadań z \textbf{Listy 6} z przedmiotu \emph{Matematyka 1}. Każde zadanie rozpoczyna się od przytoczenia pełnej treści, po czym następuje zestawienie niezbędnego aparatu pojęciowego, definicji, twierdzeń oraz kryteriów analitycznych. W dalszej części przedstawiono wyczerpujące kroki analityczne i algebraiczne wraz z uzasadnieniem każdego przejścia, a całość zwieńczona jest jednoznacznym podsumowaniem.
\end{abstract}

\vspace{0.3cm}
{
\hypersetup{linkcolor=charcoal}
\tableofcontents
}
\vspace{0.5cm}

% ------------------------------------------------------------------------------
% DOŁĄCZANIE POSZCZEGÓLNYCH ZADAŃ
% ------------------------------------------------------------------------------

\clearpage
\section{Zadanie 1: Suma współczynników dwumianowych Newtona}
\label{sec:zadanie1}

\begin{tresczadania}
Udowodnij, że dla dowolnej liczby całkowitej $N \geqslant 0$ zachodzi tożsamość:
\begin{equation*}
    \sum_{k=0}^{N} \frac{N!}{k!(N-k)!} = 2^N.
\end{equation*}
\end{tresczadania}

\vspace{0.2cm}

\begin{teoriabox}
\begin{enumerate}
    \item \textbf{Definicja symbolu Newtona (współczynnika dwumianowego):}
    Dla liczb całkowitych $N, k$ spełniających warunek $0 \leqslant k \leqslant N$, współczynnik dwumianowy $\binom{N}{k}$ definiuje się wzorem:
    \begin{equation}
        \binom{N}{k} \coloneqq \frac{N!}{k!(N-k)!},
        \label{eq:newton_def}
    \end{equation}
    z formalną konwencją $0! = 1$. Określa on liczbę kombinacji $k$-elementowych ze zbioru $N$-elementowego.

    \item \textbf{Twierdzenie o dwumianie Newtona:}
    Dla dowolnych liczb rzeczywistych $a, b \in \mathbb{R}$ oraz wykładnika $N \in \mathbb{N}_0$:
    \begin{equation}
        (a + b)^N = \sum_{k=0}^{N} \binom{N}{k} a^{N-k} b^k.
        \label{eq:binomial_theorem}
    \end{equation}

    \item \textbf{Interpretacja w teorii mnogości:}
    Liczba wszystkich możliwych podzbiorów zbioru $N$-elementowego wynosi $|\mathcal{P}(X)| = 2^N$. Ponieważ każdy podzbiór ma moc $k \in \{0, 1, \dots, N\}$, a liczba podzbiorów o ustalonej mocy $k$ wynosi dokładnie $\binom{N}{k}$, całkowita liczba podzbiorów jest równa sumie $\sum_{k=0}^N \binom{N}{k} = 2^N$.
\end{enumerate}
\end{teoriabox}

\vspace{0.2cm}

\subsection*{Rozwiązanie krok po kroku}

\subsubsection*{Metoda 1: Zastosowanie wzoru dwumianowego Newtona (analityczna)}

\begin{enumerate}[label=\textbf{Krok \arabic*.},leftmargin=*]
    \item \textbf{Identyfikacja składników sumy.}
    Zauważamy, że lewa strona tezy $L$ składa się ze składników będących z definicji~\eqref{eq:newton_def} współczynnikami dwumianowymi:
    \begin{equation*}
        L = \sum_{k=0}^{N} \frac{N!}{k!(N-k)!} = \sum_{k=0}^{N} \binom{N}{k}.
    \end{equation*}

    \item \textbf{Wybór parametrów wzoru dwumianowego.}
    Rozpatrujemy tożsamość dwumianową~\eqref{eq:binomial_theorem} dla szczególnego doboru stałych $a = 1$ oraz $b = 1$:
    \begin{equation*}
        (1 + 1)^N = \sum_{k=0}^{N} \binom{N}{k} \cdot 1^{N-k} \cdot 1^k.
    \end{equation*}

    \item \textbf{Uproszczenie potęg jedynki.}
    Dla dowolnych wykładników całkowitych potęgi liczby $1$ są równe $1$, czyli $1^{N-k} = 1$ oraz $1^k = 1$. Zatem:
    \begin{equation*}
        (1 + 1)^N = \sum_{k=0}^{N} \binom{N}{k} \cdot 1 = \sum_{k=0}^{N} \binom{N}{k}.
    \end{equation*}

    \item \textbf{Obliczenie wartości lewej strony i wniosek.}
    Ponieważ $1 + 1 = 2$, bezpośrednio otrzymujemy:
    \begin{equation*}
        2^N = \sum_{k=0}^{N} \binom{N}{k} = \sum_{k=0}^{N} \frac{N!}{k!(N-k)!}.
    \end{equation*}
    Lewa strona jest tożsamościowo równa prawej ($L = P$), co kończy dowód.
\end{enumerate}

\subsubsection*{Metoda 2: Dowód przez indukcję matematyczną (konstruktywny)}

\begin{enumerate}[label=\textbf{Krok \arabic*.},leftmargin=*]
    \item \textbf{Baza indukcji ($N = 0$).}
    Dla $N = 0$ lewa strona wynosi:
    \begin{equation*}
        L_0 = \sum_{k=0}^{0} \frac{0!}{k!(0-k)!} = \frac{0!}{0! \cdot 0!} = \frac{1}{1 \cdot 1} = 1, \qquad P_0 = 2^0 = 1.
    \end{equation*}
    Zatem $L_0 = P_0$, baza indukcji jest prawdziwa.

    \item \textbf{Założenie indukcyjne.}
    Zakładamy, że dla ustalonego $N \geqslant 0$ zachodzi: $\sum_{k=0}^{N} \binom{N}{k} = 2^N$.

    \item \textbf{Krok indukcyjny dla $N + 1$.}
    Stosujemy tożsamość trójkąta Pascala $\binom{N+1}{k} = \binom{N}{k-1} + \binom{N}{k}$ dla $1 \leqslant k \leqslant N$:
    \begin{align*}
        \sum_{k=0}^{N+1} \binom{N+1}{k} &= \binom{N+1}{0} + \sum_{k=1}^{N} \left[ \binom{N}{k-1} + \binom{N}{k} \right] + \binom{N+1}{N+1} \\
        &= 1 + \sum_{k=1}^{N} \binom{N}{k-1} + \sum_{k=1}^{N} \binom{N}{k} + 1 \\
        &= \sum_{j=0}^{N} \binom{N}{j} + \sum_{k=0}^{N} \binom{N}{k} = 2^N + 2^N = 2 \cdot 2^N = 2^{N+1}.
    \end{align*}
    Na mocy zasady indukcji matematycznej tożsamość zachodzi dla każdego $N \in \mathbb{N}_0$.
\end{enumerate}

\vspace{0.2cm}

\begin{odpowiedzbox}[Podsumowanie Zadania 1]
Tożsamość $\displaystyle\sum_{k=0}^{N} \frac{N!}{k!(N-k)!} = 2^N$ została udowodniona dwoma niezależnymi, ścisłymi metodami. Stanowi ona analityczne odzwierciedlenie faktu, że suma wszystkich podzbiorów zbioru o mocy $N$ wynosi dokładnie $2^N$.
\end{odpowiedzbox}


\clearpage
\section{Zadanie 2: Badanie rozbieżności szeregów z Listy 5}
\label{sec:zadanie2}

\begin{tresczadania}
Które z szeregów $\sum a_n$ dla ciągów $a_n$ z ćwiczenia 2 listy 5 są rozbieżne?
\begin{enumerate}[label=(\alph*),leftmargin=2em]
    \item $\displaystyle a_n = 1 + \frac{1}{n} + \frac{1}{n^2}$
    \item $\displaystyle a_n = \frac{n^2+2}{n^3-5n+1}$
    \item $\displaystyle a_n = \frac{|\cos(\pi n)|}{n}$
    \item $\displaystyle a_n = \sqrt{n^2+n} - n$
    \item $\displaystyle a_n = \frac{(n+1)! - n! + (n-1)!}{(n+1)! + n!}$
\end{enumerate}
\end{tresczadania}

\vspace{0.2cm}

\begin{teoriabox}
\begin{enumerate}
    \item \textbf{Warunek konieczny zbieżności szeregu:}
    Jeżeli szereg liczbowy $\sum_{n=1}^{\infty} a_n$ jest zbieżny, to granica jego wyrazu ogólnego wynosi zero:
    \begin{equation}
        \sum_{n=1}^{\infty} a_n < \infty \implies \lim_{n \to \infty} a_n = 0.
        \label{eq:wk_zbieznosci}
    \end{equation}

    \item \textbf{Kryterium rozbieżności (z prawa kontrapozycji):}
    Jeżeli wyraz ogólny szeregu nie zbiega do zera lub granica nie istnieje:
    \begin{equation}
        \lim_{n \to \infty} a_n \neq 0 \quad \lor \quad \lim_{n \to \infty} a_n \text{ nie istnieje} \implies \sum_{n=1}^{\infty} a_n \text{ jest \textbf{rozbieżny}}.
        \label{eq:kryt_rozbieznosci}
    \end{equation}

    \item \textbf{Kryterium ilorazowe (graniczne kryterium porównawcze):}
    Gdy $\lim_{n \to \infty} a_n = 0$, warunek konieczny nie rozstrzyga o zbieżności. Jeżeli $a_n > 0$ oraz $b_n > 0$, a ponadto istnieje właściwa i dodatnia granica:
    \begin{equation}
        k = \lim_{n \to \infty} \frac{a_n}{b_n} \in (0, \infty),
        \label{eq:kryt_ilorazowe}
    \end{equation}
    to szeregi $\sum a_n$ oraz $\sum b_n$ mają ten sam charakter zbieżności. W szczególności, gdy $b_n = \frac{1}{n}$ (szereg harmoniczny, rozbieżny), to szereg $\sum a_n$ jest również rozbieżny.
\end{enumerate}
\end{teoriabox}

\vspace{0.2cm}

\subsection*{Rozwiązania szczegółowe poszczególnych podpunktów}

% -----------------------------------------------------------------------------
\subsubsection*{Podpunkt (a): $\displaystyle a_n = 1 + \frac{1}{n} + \frac{1}{n^2}$}

\begin{enumerate}[label=\textbf{Krok \arabic*.},leftmargin=*]
    \item \textbf{Badanie granicy wyrazu ogólnego.}
    Obliczamy granicę ciągu przy $n \to \infty$:
    \begin{equation*}
        \lim_{n \to \infty} a_n = \lim_{n \to \infty} \left(1 + \frac{1}{n} + \frac{1}{n^2}\right) = 1 + 0 + 0 = 1.
    \end{equation*}
    \item \textbf{Ocena warunku koniecznego i wniosek.}
    Ponieważ $\lim_{n \to \infty} a_n = 1 \neq 0$, warunek konieczny zbieżności nie jest spełniony. Zatem szereg $\displaystyle\sum_{n=1}^{\infty} \left(1 + \frac{1}{n} + \frac{1}{n^2}\right)$ jest \textbf{rozbieżny}.
\end{enumerate}

% -----------------------------------------------------------------------------
\subsubsection*{Podpunkt (b): $\displaystyle a_n = \frac{n^2+2}{n^3-5n+1}$}

\begin{enumerate}[label=\textbf{Krok \arabic*.},leftmargin=*]
    \item \textbf{Badanie granicy wyrazu ogólnego.}
    Dzieląc licznik i mianownik przez $n^3$, otrzymujemy:
    \begin{equation*}
        \lim_{n \to \infty} a_n = \lim_{n \to \infty} \frac{\frac{1}{n} + \frac{2}{n^3}}{1 - \frac{5}{n^2} + \frac{1}{n^3}} = \frac{0}{1} = 0.
    \end{equation*}
    Warunek konieczny jest spełniony, co wymaga dalszego badania rzędu zbieżności.
    \item \textbf{Zastosowanie kryterium ilorazowego.}
    Dla dużych $n$ wyrazy są dodatnie, a asymptotycznie $a_n \sim \frac{n^2}{n^3} = \frac{1}{n}$. Przyjmijmy rozbieżny szereg harmoniczny $b_n = \frac{1}{n}$:
    \begin{align*}
        k &= \lim_{n \to \infty} \frac{a_n}{b_n} = \lim_{n \to \infty} \frac{\frac{n^2+2}{n^3-5n+1}}{\frac{1}{n}} = \lim_{n \to \infty} \frac{n(n^2+2)}{n^3-5n+1} \\
        &= \lim_{n \to \infty} \frac{n^3+2n}{n^3-5n+1} = \lim_{n \to \infty} \frac{1 + \frac{2}{n^2}}{1 - \frac{5}{n^2} + \frac{1}{n^3}} = 1.
    \end{align*}
    \item \textbf{Wniosek.}
    Ponieważ $k = 1 \in (0, \infty)$, a szereg harmoniczny $\sum \frac{1}{n}$ jest rozbieżny, badany szereg $\displaystyle\sum_{n=1}^{\infty} \frac{n^2+2}{n^3-5n+1}$ jest również \textbf{rozbieżny}.
\end{enumerate}

% -----------------------------------------------------------------------------
\subsubsection*{Podpunkt (c): $\displaystyle a_n = \frac{|\cos(\pi n)|}{n}$}

\begin{enumerate}[label=\textbf{Krok \arabic*.},leftmargin=*]
    \item \textbf{Uproszczenie licznika.}
    Dla każdego $n \in \mathbb{N}$ mamy tożsamość $\cos(\pi n) = (-1)^n$, skąd:
    \begin{equation*}
        |\cos(\pi n)| = |(-1)^n| = 1.
    \end{equation*}
    \item \textbf{Postać wyrazu ogólnego i wniosek.}
    Ciąg upraszcza się do postaci $a_n = \frac{1}{n}$. Odpowiadający mu szereg to klasyczny szereg harmoniczny:
    \begin{equation*}
        \sum_{n=1}^{\infty} a_n = \sum_{n=1}^{\infty} \frac{1}{n}.
    \end{equation*}
    Szereg harmoniczny jest powszechnie znanym szeregiem rozbieżnym do $+\infty$. Zatem badany szereg jest \textbf{rozbieżny}.
\end{enumerate}

% -----------------------------------------------------------------------------
\subsubsection*{Podpunkt (d): $\displaystyle a_n = \sqrt{n^2+n} - n$}

\begin{enumerate}[label=\textbf{Krok \arabic*.},leftmargin=*]
    \item \textbf{Mnożenie przez sprzężenie.}
    Likwidujemy symbol nieoznaczony $[\infty - \infty]$:
    \begin{align*}
        a_n &= \frac{(\sqrt{n^2+n} - n)(\sqrt{n^2+n} + n)}{\sqrt{n^2+n} + n} = \frac{(n^2+n) - n^2}{\sqrt{n^2(1 + \frac{1}{n})} + n} \\
        &= \frac{n}{n\sqrt{1 + \frac{1}{n}} + n} = \frac{1}{\sqrt{1 + \frac{1}{n}} + 1}.
    \end{align*}
    \item \textbf{Badanie granicy wyrazu ogólnego.}
    \begin{equation*}
        \lim_{n \to \infty} a_n = \lim_{n \to \infty} \frac{1}{\sqrt{1 + \frac{1}{n}} + 1} = \frac{1}{\sqrt{1+0} + 1} = \frac{1}{2}.
    \end{equation*}
    \item \textbf{Ocena i wniosek.}
    Ponieważ $\lim_{n \to \infty} a_n = \frac{1}{2} \neq 0$, szereg nie spełnia warunku koniecznego zbieżności. Stąd szereg $\displaystyle\sum_{n=1}^{\infty} \left(\sqrt{n^2+n} - n\right)$ jest \textbf{rozbieżny}.
\end{enumerate}

% -----------------------------------------------------------------------------
\subsubsection*{Podpunkt (e): $\displaystyle a_n = \frac{(n+1)! - n! + (n-1)!}{(n+1)! + n!}$}

\begin{enumerate}[label=\textbf{Krok \arabic*.},leftmargin=*]
    \item \textbf{Uproszczenie ułamka z silniami.}
    Dzielimy licznik i mianownik przez $n!$, pamiętając, że $(n+1)! = (n+1)n!$ oraz $(n-1)! = \frac{n!}{n}$:
    \begin{equation*}
        a_n = \frac{(n+1) - 1 + \frac{1}{n}}{(n+1) + 1} = \frac{n + \frac{1}{n}}{n + 2} = \frac{n^2 + 1}{n^2 + 2n}.
    \end{equation*}
    \item \textbf{Badanie granicy wyrazu ogólnego.}
    \begin{equation*}
        \lim_{n \to \infty} a_n = \lim_{n \to \infty} \frac{1 + \frac{1}{n^2}}{1 + \frac{2}{n}} = \frac{1+0}{1+0} = 1.
    \end{equation*}
    \item \textbf{Ocena i wniosek.}
    Ponieważ $\lim_{n \to \infty} a_n = 1 \neq 0$, warunek konieczny nie jest spełniony. Zatem szereg $\displaystyle\sum_{n=1}^{\infty} \frac{(n+1)! - n! + (n-1)!}{(n+1)! + n!}$ jest \textbf{rozbieżny}.
\end{enumerate}

\vspace{0.2cm}

\begin{odpowiedzbox}[Podsumowanie i odpowiedź do Zadania 2]
\begin{center}
\small
\resizebox{0.96\linewidth}{!}{%
\begin{tabular}{c l c c l}
\toprule
\textbf{Podpunkt} & \textbf{Ciąg $a_n$} & $\mathbf{\lim a_n}$ & \textbf{Warunek konieczny} & \textbf{Werdykt i uzasadnienie} \\
\midrule
(a) & $1 + \frac{1}{n} + \frac{1}{n^2}$ & $1$ & Nie spełniony & \textbf{Rozbieżny} ($\lim a_n = 1 \neq 0$) \\
(b) & $\frac{n^2+2}{n^3-5n+1}$ & $0$ & Spełniony & \textbf{Rozbieżny} (ilorazowe z $\sum \frac{1}{n}$) \\
(c) & $\frac{|\cos(\pi n)|}{n} = \frac{1}{n}$ & $0$ & Spełniony & \textbf{Rozbieżny} (szereg harmoniczny) \\
(d) & $\sqrt{n^2+n} - n$ & $\frac{1}{2}$ & Nie spełniony & \textbf{Rozbieżny} ($\lim a_n = \frac{1}{2} \neq 0$) \\
(e) & $\frac{(n+1)! - n! + (n-1)!}{(n+1)! + n!}$ & $1$ & Nie spełniony & \textbf{Rozbieżny} ($\lim a_n = 1 \neq 0$) \\
\bottomrule
\end{tabular}%
}
\end{center}

\vspace{0.2cm}
\noindent\textbf{Ostateczna konkluzja:}\\
\textbf{Wszystkie} badane szeregi (a), (b), (c), (d) oraz (e) są \textbf{rozbieżne}. Szeregi (a), (d) i (e) rozbiegają się wskutek niespełnienia warunku koniecznego zbieżności, natomiast szeregi (b) i (c) rozbiegają się ze względu na zbyt powolny zanik wyrazu ogólnego (rząd harmoniczny $\sim \frac{1}{n}$).
\end{odpowiedzbox}


\clearpage
\section{Zadanie 3: Badanie zbieżności szeregów za pomocą kryteriów}
\label{sec:zadanie3}

\begin{tresczadania}
Zbadaj zbieżność następujących szeregów (korzystając z kryteriów zbieżności):
\begin{enumerate}[label=(\alph*),leftmargin=2em]
    \item $\displaystyle \sum_{n=1}^{\infty} \frac{2^n}{n!}$
    \item $\displaystyle \sum_{n=1}^{\infty} \frac{3^n}{n^2}$
    \item $\displaystyle 1 + \sum_{n=2}^{\infty} \frac{1}{n-1}$
    \item $\displaystyle \sum_{n=1}^{\infty} \left(\frac{\pi}{2e}\right)^n$
    \item $\displaystyle \sum_{n=1}^{\infty} \frac{7^{n+2}}{5^n}$
    \item $\displaystyle \sum_{n=1}^{\infty} \left(\frac{1}{2^n} - \frac{1}{2^{n+1}}\right)$
\end{enumerate}
\end{tresczadania}

\vspace{0.2cm}

\begin{teoriabox}
\begin{enumerate}
    \item \textbf{Kryterium d'Alemberta (ilorazowe):}
    Dla szeregu o wyrazach dodatnich ($a_n > 0$), jeżeli istnieje granica:
    \begin{equation}
        g = \lim_{n \to \infty} \frac{a_{n+1}}{a_n},
        \label{eq:dalembert_crit}
    \end{equation}
    to:
    \begin{itemize}
        \item gdy $g < 1$, szereg jest \textbf{zbieżny},
        \item gdy $g > 1$, szereg jest \textbf{rozbieżny},
        \item gdy $g = 1$, kryterium nie rozstrzyga o zbieżności.
    \end{itemize}

    \item \textbf{Kryterium Cauchy'ego (pierwiastkowe):}
    Dla szeregu o wyrazach nieujemnych ($a_n \geqslant 0$), jeżeli istnieje granica:
    \begin{equation}
        g = \lim_{n \to \infty} \sqrt[n]{a_n},
        \label{eq:cauchy_crit}
    \end{equation}
    to gdy $g < 1$ szereg jest \textbf{zbieżny}, gdy $g > 1$ jest \textbf{rozbieżny}, a gdy $g = 1$ kryterium nie rozstrzyga.

    \item \textbf{Zbieżność szeregu geometrycznego:}
    Szereg $\sum_{n=1}^{\infty} c \cdot q^n$ ($c \neq 0$) jest \textbf{zbieżny} $\iff |q| < 1$. Wtedy jego suma wynosi $S = \frac{a_1}{1 - q}$. Dla $|q| \geqslant 1$ szereg jest \textbf{rozbieżny}.

    \item \textbf{Szereg harmoniczny rzędu $\alpha$:}
    Szereg $\sum_{n=1}^{\infty} \frac{1}{n^\alpha}$ jest zbieżny $\iff \alpha > 1$. Dla $\alpha \leqslant 1$ szereg jest rozbieżny.

    \item \textbf{Niezmienniczość zbieżności ze względu na początkowe wyrazy:}
    Modyfikacja skończonej liczby wyrazów (np. dodanie stałej na początku) nie wpływa na fakt zbieżności bądź rozbieżności szeregu.
\end{enumerate}
\end{teoriabox}

\vspace{0.2cm}

\subsection*{Rozwiązania szczegółowe poszczególnych podpunktów}

% -----------------------------------------------------------------------------
\subsubsection*{Podpunkt (a): $\displaystyle \sum_{n=1}^{\infty} \frac{2^n}{n!}$}

\begin{enumerate}[label=\textbf{Krok \arabic*.},leftmargin=*]
    \item \textbf{Wybór metody.}
    Mamy $a_n = \frac{2^n}{n!} > 0$. Obecność silni w mianowniku wskazuje na celowość zastosowania \textbf{kryterium d'Alemberta}.
    \item \textbf{Obliczenie ilorazu $\frac{a_{n+1}}{a_n}$.}
    \begin{equation*}
        \frac{a_{n+1}}{a_n} = \frac{2^{n+1}}{(n+1)!} \cdot \frac{n!}{2^n} = \frac{2 \cdot 2^n \cdot n!}{(n+1)n! \cdot 2^n} = \frac{2}{n+1}.
    \end{equation*}
    \item \textbf{Granica ilorazu i wniosek.}
    \begin{equation*}
        g = \lim_{n \to \infty} \frac{a_{n+1}}{a_n} = \lim_{n \to \infty} \frac{2}{n+1} = 0.
    \end{equation*}
    Ponieważ $g = 0 < 1$, na mocy kryterium d'Alemberta szereg $\displaystyle\sum_{n=1}^{\infty} \frac{2^n}{n!}$ jest \textbf{zbieżny}.
\end{enumerate}

% -----------------------------------------------------------------------------
\subsubsection*{Podpunkt (b): $\displaystyle \sum_{n=1}^{\infty} \frac{3^n}{n^2}$}

\begin{enumerate}[label=\textbf{Krok \arabic*.},leftmargin=*]
    \item \textbf{Zastosowanie kryterium d'Alemberta.}
    Dla $a_n = \frac{3^n}{n^2} > 0$ wyznaczamy iloraz kolejnych wyrazów:
    \begin{equation*}
        \frac{a_{n+1}}{a_n} = \frac{3^{n+1}}{(n+1)^2} \cdot \frac{n^2}{3^n} = 3 \cdot \left(\frac{n}{n+1}\right)^2 = 3 \cdot \left(\frac{1}{1 + \frac{1}{n}}\right)^2.
    \end{equation*}
    \item \textbf{Granica ilorazu i wniosek.}
    \begin{equation*}
        g = \lim_{n \to \infty} \frac{a_{n+1}}{a_n} = 3 \cdot 1^2 = 3.
    \end{equation*}
    Ponieważ $g = 3 > 1$, na mocy kryterium d'Alemberta szereg $\displaystyle\sum_{n=1}^{\infty} \frac{3^n}{n^2}$ jest \textbf{rozbieżny} (co potwierdza również warunek konieczny: $\lim a_n = +\infty \neq 0$).
\end{enumerate}

% -----------------------------------------------------------------------------
\subsubsection*{Podpunkt (c): $\displaystyle 1 + \sum_{n=2}^{\infty} \frac{1}{n-1}$}

\begin{enumerate}[label=\textbf{Krok \arabic*.},leftmargin=*]
    \item \textbf{Przesunięcie wskaźnika sumowania.}
    Wprowadzamy nową zmienną sumowania $k = n - 1$. Dla $n = 2$ mamy $k = 1$, zatem:
    \begin{equation*}
        \sum_{n=2}^{\infty} \frac{1}{n-1} = \sum_{k=1}^{\infty} \frac{1}{k}.
    \end{equation*}
    \item \textbf{Rozpoznanie szeregu i wniosek.}
    Jest to klasyczny szereg harmoniczny rzędu $\alpha = 1$, o którym wiadomo, że jest rozbieżny do $+\infty$. Dodanie stałej $1$ nie wpływa na zbieżność:
    \begin{equation*}
        1 + \sum_{n=2}^{\infty} \frac{1}{n-1} = 1 + \infty = \infty.
    \end{equation*}
    Badany szereg jest \textbf{rozbieżny}.
\end{enumerate}

% -----------------------------------------------------------------------------
\subsubsection*{Podpunkt (d): $\displaystyle \sum_{n=1}^{\infty} \left(\frac{\pi}{2e}\right)^n$}

\begin{enumerate}[label=\textbf{Krok \arabic*.},leftmargin=*]
    \item \textbf{Identyfikacja ilorazu szeregu geometrycznego.}
    Mamy szereg geometryczny o ilorazie $q = \frac{\pi}{2e}$.
    \item \textbf{Oszacowanie wartości $q$.}
    Przyjmując znane przybliżenia $\pi \approx 3{,}14159$ oraz $e \approx 2{,}71828 \implies 2e \approx 5{,}43656$, widzimy, że $\pi < 2e$, skąd:
    \begin{equation*}
        0 < q = \frac{\pi}{2e} \approx 0{,}5779 < 1.
    \end{equation*}
    \item \textbf{Wniosek.}
    Ponieważ $|q| < 1$, na mocy twierdzenia o zbieżności szeregu geometrycznego (lub kryterium Cauchy'ego: $\lim \sqrt[n]{a_n} = \frac{\pi}{2e} < 1$), szereg $\displaystyle\sum_{n=1}^{\infty} \left(\frac{\pi}{2e}\right)^n$ jest \textbf{zbieżny}.
\end{enumerate}

% -----------------------------------------------------------------------------
\subsubsection*{Podpunkt (e): $\displaystyle \sum_{n=1}^{\infty} \frac{7^{n+2}}{5^n}$}

\begin{enumerate}[label=\textbf{Krok \arabic*.},leftmargin=*]
    \item \textbf{Przekształcenie algebraiczne.}
    \begin{equation*}
        a_n = \frac{7^{n+2}}{5^n} = \frac{7^2 \cdot 7^n}{5^n} = 49 \cdot \left(\frac{7}{5}\right)^n.
    \end{equation*}
    \item \textbf{Zastosowanie kryterium Cauchy'ego.}
    \begin{equation*}
        g = \lim_{n \to \infty} \sqrt[n]{a_n} = \lim_{n \to \infty} \sqrt[n]{49} \cdot \frac{7}{5} = 1 \cdot \frac{7}{5} = \frac{7}{5} = 1{,}4.
    \end{equation*}
    \item \textbf{Wniosek.}
    Ponieważ $g = 1{,}4 > 1$, szereg $\displaystyle\sum_{n=1}^{\infty} \frac{7^{n+2}}{5^n}$ jest \textbf{rozbieżny} (ponadto $\lim a_n = +\infty \neq 0$).
\end{enumerate}

% -----------------------------------------------------------------------------
\subsubsection*{Podpunkt (f): $\displaystyle \sum_{n=1}^{\infty} \left(\frac{1}{2^n} - \frac{1}{2^{n+1}}\right)$}

\begin{enumerate}[label=\textbf{Krok \arabic*.},leftmargin=*]
    \item \textbf{Metoda teleskopowa.}
    Dla $b_n = \frac{1}{2^n}$ wyraz ogólny ma postać $a_n = b_n - b_{n+1}$. Suma częściowa $S_k$ wynosi:
    \begin{equation*}
        S_k = \sum_{n=1}^{k} (b_n - b_{n+1}) = b_1 - b_{k+1} = \frac{1}{2} - \frac{1}{2^{k+1}}.
    \end{equation*}
    Przechodząc do granicy przy $k \to \infty$:
    \begin{equation*}
        S = \lim_{k \to \infty} S_k = \lim_{k \to \infty} \left(\frac{1}{2} - \frac{1}{2^{k+1}}\right) = \frac{1}{2} - 0 = \frac{1}{2}.
    \end{equation*}
    \item \textbf{Metoda geometryczna (równoważna).}
    Upraszczając ułamek: $a_n = \frac{2 - 1}{2^{n+1}} = \frac{1}{2^{n+1}} = \frac{1}{4} \cdot \left(\frac{1}{2}\right)^{n-1}$. Jest to szereg geometryczny o ilorazie $q = \frac{1}{2} < 1$, którego suma wynosi:
    \begin{equation*}
        S = \frac{a_1}{1 - q} = \frac{\frac{1}{4}}{1 - \frac{1}{2}} = \frac{1}{2}.
    \end{equation*}
    \item \textbf{Wniosek.}
    Szereg $\displaystyle\sum_{n=1}^{\infty} \left(\frac{1}{2^n} - \frac{1}{2^{n+1}}\right)$ jest \textbf{zbieżny}, a jego suma wynosi dokładnie $\frac{1}{2}$.
\end{enumerate}

\vspace{0.2cm}

\begin{odpowiedzbox}[Podsumowanie wyników Zadania 3]
\begin{center}
\small
\resizebox{0.96\linewidth}{!}{%
\begin{tabular}{c l c l}
\toprule
\textbf{Podpunkt} & \textbf{Szereg} & \textbf{Werdykt} & \textbf{Zastosowane kryterium / Uzasadnienie} \\
\midrule
(a) & $\sum_{n=1}^{\infty} \frac{2^n}{n!}$ & \textbf{Zbieżny} & Kryterium d'Alemberta ($g = 0 < 1$) \\
(b) & $\sum_{n=1}^{\infty} \frac{3^n}{n^2}$ & \textbf{Rozbieżny} & Kryterium d'Alemberta ($g = 3 > 1$) \\
(c) & $1 + \sum_{n=2}^{\infty} \frac{1}{n-1}$ & \textbf{Rozbieżny} & Szereg harmoniczny $\sum \frac{1}{k}$ ($\alpha = 1$) \\
(d) & $\sum_{n=1}^{\infty} \left(\frac{\pi}{2e}\right)^n$ & \textbf{Zbieżny} & Szereg geometryczny ($|q| \approx 0{,}578 < 1$) \\
(e) & $\sum_{n=1}^{\infty} \frac{7^{n+2}}{5^n}$ & \textbf{Rozbieżny} & Szereg geometryczny ($q = 1{,}4 > 1$) / Cauchy \\
(f) & $\sum_{n=1}^{\infty} \left(\frac{1}{2^n} - \frac{1}{2^{n+1}}\right)$ & \textbf{Zbieżny} & Szereg teleskopowy ($S = \frac{1}{2}$) \\
\bottomrule
\end{tabular}%
}
\end{center}
\end{odpowiedzbox}


\clearpage
\section{Zadanie 4: Obliczanie n-tej sumy i badanie zbieżności szeregów}
\label{sec:zadanie4}

\begin{tresczadania}
Obliczyć $n$-tą sumę i zbadać zbieżność szeregów:
\begin{enumerate}[label=(\alph*),leftmargin=2em]
    \item $\displaystyle \sum_{n=1}^{\infty} \ln\left(1 + \frac{1}{n}\right)$
    \item $\displaystyle 1 + \sum_{n=2}^{\infty} \ln\left(1 - \frac{1}{n^2}\right)$
\end{enumerate}
\end{tresczadania}

\vspace{0.2cm}

\begin{teoriabox}
\begin{enumerate}
    \item \textbf{Ciąg sum częściowych ($n$-ta suma):}
    Dla szeregu liczb rzeczywistych $\sum_{k=1}^{\infty} a_k$, ciąg sum częściowych definiuje się jako:
    \begin{equation}
        S_n \coloneqq \sum_{k=1}^{n} a_k = a_1 + a_2 + \dots + a_n.
        \label{eq:partial_sum_def}
    \end{equation}

    \item \textbf{Definicja zbieżności szeregu:}
    Szereg jest \textbf{zbieżny} wtedy i tylko wtedy, gdy istnieje skończona granica ciągu sum częściowych:
    \begin{equation*}
        S = \lim_{n \to \infty} S_n \in \mathbb{R}.
    \end{equation*}
    Wartość $S$ nazywamy \textbf{sumą szeregu}. Gdy granica jest nieskończona lub nie istnieje, szereg jest \textbf{rozbieżny}.

    \item \textbf{Własności logarytmu naturalnego i iloczynów:}
    Dla liczb dodatnich zachodzą prawa:
    \begin{equation*}
        \ln(x \cdot y) = \ln x + \ln y, \qquad \ln\left(\frac{x}{y}\right) = \ln x - \ln y \implies \sum_{k=1}^{n} \ln(u_k) = \ln\left( \prod_{k=1}^{n} u_k \right).
    \end{equation*}

    \item \textbf{Ciągłość funkcji logarytmicznej:}
    Z uwagi na ciągłość funkcji $\ln$ na $(0, \infty)$, przejście graniczne można przenieść do argumentu logarytmu: $\lim \ln(u_n) = \ln(\lim u_n)$, o ile $\lim u_n > 0$.
\end{enumerate}
\end{teoriabox}

\vspace{0.2cm}

\subsection*{Rozwiązania szczegółowe}

% -----------------------------------------------------------------------------
\subsubsection*{Podpunkt (a): $\displaystyle \sum_{n=1}^{\infty} \ln\left(1 + \frac{1}{n}\right)$}

\begin{enumerate}[label=\textbf{Krok \arabic*.},leftmargin=*]
    \item \textbf{Przekształcenie wyrazu ogólnego.}
    Zapisujemy argument pod wspólnym mianownikiem i rozbijamy logarytm ilorazu:
    \begin{equation*}
        a_k = \ln\left(1 + \frac{1}{k}\right) = \ln\left(\frac{k+1}{k}\right) = \ln(k+1) - \ln k.
    \end{equation*}

    \item \textbf{Wyznaczenie $n$-tej sumy częściowej $S_n$.}
    Zapisujemy sumę wyrazów od $k = 1$ do $n$:
    \begin{align*}
        S_n &= \sum_{k=1}^{n} \big( \ln(k+1) - \ln k \big) \\
        &= (\ln 2 - \ln 1) + (\ln 3 - \ln 2) + (\ln 4 - \ln 3) + \dots + \big( \ln(n+1) - \ln n \big).
    \end{align*}
    Wyrazy pośrednie redukują się wzajemnie (zwijanie teleskopowe):
    \begin{equation*}
        S_n = \ln(n+1) - \ln 1 = \ln(n+1) - 0 = \ln(n+1).
    \end{equation*}
    (Iloczynowo: $S_n = \ln\left(\prod_{k=1}^n \frac{k+1}{k}\right) = \ln\left(\frac{2}{1} \cdot \frac{3}{2} \cdots \frac{n+1}{n}\right) = \ln(n+1)$).

    \item \textbf{Badanie zbieżności szeregu.}
    Obliczamy granicę sum częściowych przy $n \to \infty$:
    \begin{equation*}
        \lim_{n \to \infty} S_n = \lim_{n \to \infty} \ln(n+1) = +\infty.
    \end{equation*}

    \item \textbf{Wniosek.}
    Ponieważ granica ciągu sum częściowych wynosi $+\infty$, szereg $\displaystyle\sum_{n=1}^{\infty} \ln\left(1 + \frac{1}{n}\right)$ jest \textbf{rozbieżny}.
\end{enumerate}

% -----------------------------------------------------------------------------
\subsubsection*{Podpunkt (b): $\displaystyle 1 + \sum_{n=2}^{\infty} \ln\left(1 - \frac{1}{n^2}\right)$}

\begin{enumerate}[label=\textbf{Krok \arabic*.},leftmargin=*]
    \item \textbf{Faktoryzacja argumentu logarytmu.}
    Dla $k \geqslant 2$ stosujemy wzór na różnicę kwadratów:
    \begin{equation*}
        1 - \frac{1}{k^2} = \frac{k^2 - 1}{k^2} = \frac{(k-1)(k+1)}{k^2} = \left(\frac{k-1}{k}\right) \cdot \left(\frac{k+1}{k}\right).
    \end{equation*}

    \item \textbf{Obliczenie sumy składników szeregu (od $k = 2$ do $n$).}
    Oznaczmy sumę częściową samej sumy szeregu przez $T_n$:
    \begin{equation*}
        T_n = \sum_{k=2}^{n} \ln\left(1 - \frac{1}{k^2}\right) = \ln\left( \prod_{k=2}^{n} \frac{(k-1)(k+1)}{k^2} \right).
    \end{equation*}
    Rozbijamy iloczyn na dwa niezależne iloczyny teleskopowe $P_n$ i $Q_n$:
    \begin{align*}
        P_n &= \prod_{k=2}^{n} \frac{k-1}{k} = \frac{1}{2} \cdot \frac{2}{3} \cdot \frac{3}{4} \cdots \frac{n-1}{n} = \frac{1}{n}, \\
        Q_n &= \prod_{k=2}^{n} \frac{k+1}{k} = \frac{3}{2} \cdot \frac{4}{3} \cdot \frac{5}{4} \cdots \frac{n+1}{n} = \frac{n+1}{2}.
    \end{align*}
    Iloczyn całkowity wynosi zatem:
    \begin{equation*}
        \prod_{k=2}^{n} \frac{(k-1)(k+1)}{k^2} = P_n \cdot Q_n = \frac{1}{n} \cdot \frac{n+1}{2} = \frac{n+1}{2n}.
    \end{equation*}
    Stąd:
    \begin{equation*}
        T_n = \ln\left(\frac{n+1}{2n}\right).
    \end{equation*}

    \item \textbf{Wyznaczenie $n$-tej sumy całego wyrażenia.}
    Całe zadane wyrażenie zawiera początkowy wyraz $1$, stąd pełna $n$-ta suma wynosi:
    \begin{equation*}
        S_n = 1 + T_n = 1 + \ln\left(\frac{n+1}{2n}\right).
    \end{equation*}

    \item \textbf{Badanie zbieżności i obliczenie sumy szeregu.}
    Obliczamy granicę przy $n \to \infty$:
    \begin{equation*}
        \lim_{n \to \infty} \frac{n+1}{2n} = \lim_{n \to \infty} \left(\frac{1}{2} + \frac{1}{2n}\right) = \frac{1}{2}.
    \end{equation*}
    Z ciągłości logarytmu:
    \begin{equation*}
        \lim_{n \to \infty} \ln\left(\frac{n+1}{2n}\right) = \ln\left(\frac{1}{2}\right) = -\ln 2.
    \end{equation*}
    Zatem granica sumy wynosi:
    \begin{equation*}
        S = \lim_{n \to \infty} S_n = 1 + \ln\left(\frac{1}{2}\right) = 1 - \ln 2.
    \end{equation*}

    \item \textbf{Wniosek.}
    Granica $S$ istnieje i jest skończona ($S = 1 - \ln 2 \approx 0{,}3069$). Zatem szereg jest \textbf{zbieżny}, a jego suma wynosi $1 - \ln 2$.
\end{enumerate}

\vspace{0.2cm}

\begin{odpowiedzbox}[Podsumowanie wyników Zadania 4]
\begin{enumerate}[label=(\alph*),leftmargin=2em]
    \item Szereg $\displaystyle\sum_{n=1}^{\infty} \ln\left(1 + \frac{1}{n}\right)$:
    \begin{itemize}
        \item $n$-ta suma częściowa: $\mathbf{S_n = \ln(n+1)}$,
        \item Zbieżność: szereg jest \textbf{rozbieżny} ($\lim_{n\to\infty} S_n = +\infty$).
    \end{itemize}
    \item Wyrażenie $\displaystyle 1 + \sum_{n=2}^{\infty} \ln\left(1 - \frac{1}{n^2}\right)$:
    \begin{itemize}
        \item $n$-ta suma: $\mathbf{S_n = 1 + \ln\left(\frac{n+1}{2n}\right)}$,
        \item Zbieżność: szereg jest \textbf{zbieżny}, a jego suma wynosi $\mathbf{S = 1 - \ln 2}$.
    \end{itemize}
\end{enumerate}
\end{odpowiedzbox}


\clearpage
\section{Zadanie 5: Zbieżność szeregu potęgowego dla dowolnego x}
\label{sec:zadanie5}

\begin{tresczadania}
Wykaż, że szereg:
\begin{equation*}
    \sum_{n=1}^{\infty} \frac{x^n}{n!}
\end{equation*}
jest zbieżny dla dowolnego $x \in \mathbb{R}$.
\end{tresczadania}

\vspace{0.2cm}

\begin{teoriabox}
\begin{enumerate}
    \item \textbf{Zbieżność bezwzględna a zbieżność zwykła:}
    Szereg $\sum_{n=1}^{\infty} a_n$ nazywamy bezwzględnie zbieżnym, jeżeli zbieżny jest szereg modułów:
    \begin{equation}
        \sum_{n=1}^{\infty} |a_n| < \infty.
        \label{eq:abs_conv_def}
    \end{equation}
    \textbf{Twierdzenie o zbieżności bezwzględnej:} Każdy szereg bezwzględnie zbieżny jest zbieżny:
    \begin{equation*}
        \sum_{n=1}^{\infty} |a_n| < \infty \implies \sum_{n=1}^{\infty} a_n \text{ jest zbieżny}.
    \end{equation*}

    \item \textbf{Kryterium d'Alemberta (wersja dla zbieżności bezwzględnej):}
    Dla szeregu o wyrazach niezerowych ($a_n \neq 0$), jeżeli istnieje granica:
    \begin{equation}
        g = \lim_{n \to \infty} \left| \frac{a_{n+1}}{a_n} \right| < 1,
        \label{eq:dalembert_modul}
    \end{equation}
    to szereg $\sum a_n$ jest bezwzględnie zbieżny, a w konsekwencji zbieżny.

    \item \textbf{Promień zbieżności szeregu potęgowego:}
    Dla szeregu potęgowego $\sum c_n x^n$ o współczynnikach $c_n = \frac{1}{n!}$, promień zbieżności wynosi:
    \begin{equation*}
        R = \lim_{n \to \infty} \left|\frac{c_n}{c_{n+1}}\right| = \lim_{n \to \infty} \frac{(n+1)!}{n!} = \lim_{n \to \infty} (n+1) = \infty.
    \end{equation*}
    Oznacza to zbieżność na całej prostej rzeczywistej $\mathbb{R} = (-\infty, +\infty)$.
\end{enumerate}
\end{teoriabox}

\vspace{0.2cm}

\subsection*{Kolejne kroki formalnego dowodu}

Ustalmy dowolną liczbę rzeczywistą $x \in \mathbb{R}$. Rozważamy wyraz ogólny szeregu:
\begin{equation*}
    a_n = \frac{x^n}{n!}, \quad n \in \mathbb{N}.
\end{equation*}

\begin{enumerate}[label=\textbf{Krok \arabic*.},leftmargin=*]
    \item \textbf{Przypadek 1: $x = 0$.}
    Dla $x = 0$ każdy wyraz ciągu wynosi:
    \begin{equation*}
        a_n = \frac{0^n}{n!} = 0 \quad \text{dla wszystkich } n \geqslant 1.
    \end{equation*}
    Suma częściowa $S_k = \sum_{n=1}^k 0 = 0$ jest ciągiem tożsamościowo stałym, którego granica wynosi $\lim S_k = 0$. Szereg jest w sposób oczywisty zbieżny.

    \item \textbf{Przypadek 2: $x \neq 0$ (badanie zbieżności bezwzględnej).}
    Dla $x \neq 0$ badamy szereg wartości bezwzględnych:
    \begin{equation*}
        \sum_{n=1}^{\infty} |a_n| = \sum_{n=1}^{\infty} \left| \frac{x^n}{n!} \right| = \sum_{n=1}^{\infty} \frac{|x|^n}{n!}.
    \end{equation*}
    Ponieważ $|x| > 0$, wszystkie wyrazy $b_n = \frac{|x|^n}{n!}$ są ściśle dodatnie ($b_n > 0$).

    \item \textbf{Zastosowanie kryterium d'Alemberta.}
    Badamy iloraz dwóch kolejnych wyrazów:
    \begin{equation*}
        \frac{b_{n+1}}{b_n} = \frac{\frac{|x|^{n+1}}{(n+1)!}}{\frac{|x|^n}{n!}} = \frac{|x|^{n+1}}{|x|^n} \cdot \frac{n!}{(n+1)!} = |x| \cdot \frac{n!}{(n+1) \cdot n!} = \frac{|x|}{n+1}.
    \end{equation*}

    \item \textbf{Obliczenie granicy ilorazu przy $n \to \infty$.}
    Liczba $|x|$ jest ustaloną stałą (nie zależy od wskaźnika $n$). Zatem:
    \begin{equation*}
        g = \lim_{n \to \infty} \frac{b_{n+1}}{b_n} = \lim_{n \to \infty} \frac{|x|}{n+1} = |x| \cdot \lim_{n \to \infty} \frac{1}{n+1} = |x| \cdot 0 = 0.
    \end{equation*}

    \item \textbf{Wnioskowanie o zbieżności bezwzględnej.}
    Niezależnie od wyboru wartości $x \neq 0$ granica wynosi:
    \begin{equation*}
        g = 0 < 1.
    \end{equation*}
    Z kryterium d'Alemberta wynika, że szereg modułów $\sum_{n=1}^{\infty} \frac{|x|^n}{n!}$ jest zbieżny dla każdego $x \neq 0$.

    \item \textbf{Wniosek ostateczny.}
    Ponieważ zbieżność bezwzględna pociąga za sobą zbieżność zwykłą, a przypadek $x=0$ jest trywialnie zbieżny, szereg $\displaystyle\sum_{n=1}^{\infty} \frac{x^n}{n!}$ jest \textbf{zbieżny dla dowolnego $x \in \mathbb{R}$}.
\end{enumerate}

\vspace{0.2cm}

\begin{odpowiedzbox}[Podsumowanie Zadania 5]
Wykazano, że szereg $\displaystyle\sum_{n=1}^{\infty} \frac{x^n}{n!}$ jest bezwzględnie zbieżny (a więc i zbieżny) na całym ciele liczb rzeczywistych $\mathbb{R}$. Silnia w mianowniku rośnie asymptotycznie szybciej niż dowolna funkcja wykładnicza $|x|^n$, co gwarantuje nieskończony promień zbieżności $R = \infty$. Szereg ten stanowi rozwinięcie funkcji wykładniczej pomniejszonej o jedynkę:
\begin{equation*}
    \sum_{n=1}^{\infty} \frac{x^n}{n!} = e^x - 1 \quad \text{dla każdego } x \in \mathbb{R}.
\end{equation*}
\end{odpowiedzbox}


\end{document}
